% Explanation certificate for heart-s5-c1 (genome 4f9f15fb-446b-4560-9704-c7b1a83617b5).
% Generated by explaneat export (schema 1.4.0) from the JSON record alone;
% content sha256 f3feba786432f89359505ee8928bcbe868228dc48e57707b5b20ab57f49b46fb.
% Drafter recorded in the JSON: ai_agent (claude). Do not edit by hand: re-render with
%   uv run python -m explaneat export render --from heart-s5-c1.json
% TikZ modules are input from \explaneatcertdir{<label>} (empty by default: same
% directory; index.tex sets it to the genome's subdirectory).
\providecommand{\explaneatcertdir}[1]{}
\begin{drafted}[appendix: heart-s5-c1 explanation certificate]
\section{Explanation certificate: heart-s5-c1}\label{app:heart-s5-c1:sec}
This section is the explanation certificate of the heart-s5-c1 network. It is rendered from the canonical JSON record \texttt{heart-s5-c1.json} (schema 1.4.0, SHA-256 \texttt{f3feba786432f893}\ldots), which also holds every evidence payload, the full graphs with weights, and the explanation record. \Cref{app:heart-s5-c1:summary} summarises the network and its coverage.

\begin{table}[htbp]
\centering\small
\begin{tabularx}{\linewidth}{@{}lX@{}}
\toprule
Dataset & Heart Disease (binary, prepared) \\
Experiment & 10seed-heart-s5, seed 5 \\
Genome & \texttt{4f9f15fb-446b-4560-9704-c7b1a83617b5} (NEAT key 9259, generation 39) \\
Explanation & \texttt{39c7fe21-518b-4517-8765-7e6008dc7d98} (36 operations) \\
Train AUC / accuracy & 0.9418 / 0.8719 (n = 242) \\
Test AUC / accuracy & 0.8869 / 0.8033 (n = 61) \\
Output activation & sigmoid \\
Phenotype & 30 nodes (26 inputs, 3 hidden, 1 output), 31 enabled edges \\
Explained model state & 36 nodes, 33 enabled edges \\
Refactoring operations & add\_\allowbreak{}identity\_\allowbreak{}node 2, split\_\allowbreak{}node 4 \\
Function preservation & max $|\Delta\,\mathrm{output}|$ = $2.22\times 10^{-16}$ on the training split \\
Annotations & 4 (3 leaf, 1 composition) \\
Structural coverage & 1 (35 of 35 candidate nodes; 4 Null) \\
Compositional coverage & 1 \\
Uncovered nodes / edges & 0 / 0 \\
Missing compositions & 0 \\
Full coverage & yes \\
Drafted by & AI-assisted draft \\
Verified by & unverified \\
Record hash (SHA-256) & \texttt{f3feba786432f893}\ldots \\
\bottomrule
\end{tabularx}
\caption{Explanation certificate summary: heart-s5-c1.}
\label{app:heart-s5-c1:summary}
\end{table}

\paragraph{Drafter's summary.} Mostly logistic regression: 19 direct edges (logit SD 3.14; train AUC 0.934 alone) plus one serial ischaemia chain (logit SD 0.67) to 0.942. An ischaemia gate ReLU(0.25 + 2.59 oldpeak + 0.87 restecg:0 + 1.15 thal:7) opens above an ST-depression threshold set by a normal resting ECG and a reversible defect and adds risk directly; it feeds a gate that passes the score on minus a resting-blood-pressure offset (more risk when BP is low), which in turn switches off a small extra age term. Oldpeak and restecg:0 act only through this chain. Same design as the 10-seed heart-s5 record (generation 37, key 9330), with three hidden nodes instead of six.

\paragraph{What resisted explanation.} Nothing resisted coverage. The serial chain 204 -$>$ 290 -$>$ 1987 (each node also read out to the output) was made into three leaves by intercepting the two chain links with identity nodes (identity\_\allowbreak{}290, identity\_\allowbreak{}1987) that become entries of the downstream leaves; the root composes them. Counter-directional signs reported, not explained: restecg:0 (normal ECG) raises the ischaemia score, trestbps enters the risk-raising BP gate negatively (low resting BP opens it), fbs:0 +0.28 directly. The age gate barely matters (readout ablation -0.001 AUC).

\subsection{Plain-language summary}\label{app:heart-s5-c1:plain}
\noindent{\small\textit{AI-assisted draft for review (unverified). Each sentence cites the claims (C01, C02, \ldots) that back it; every cited claim was computed on the model and the training split and passed its automatic checks. Magnitudes are changes in predicted risk for a typical patient; confidence words follow the test-split check and n. The text describes what the model does, not the world.}}\par\smallskip
\paragraph{Summary (drafted).} For a typical patient (a 56-year-old man with no chest pain symptoms and median test results) the model predicts a risk of heart disease of about 22\%; it adds up most of its inputs independently, and a chain of three parts adds to that sum: a signal from ST depression on exercise, the thallium scan and the resting ECG that acts on 116 of 242 training patients, a blood pressure step that reads that signal, and an extra age term that acts on 85 patients.\,{\scriptsize\textcolor{NavyBlue}{[C32, C40, C51]}} These add up independently, and these lower the risk: being a woman takes a typical patient's predicted risk from about 22\% to about 3\% and fasting blood sugar above 120 mg/dL to about 11\%; a higher maximum heart rate lowers it steadily, from about 27\% at 118 to about 18\% at 175 beats per minute; and through this sum alone typical angina lowers it to about 2\%, non-anginal chest pain to about 5\% and atypical angina to about 11\% (chest pain type also acts through the blood pressure step, below).\,{\scriptsize\textcolor{NavyBlue}{[C12, C17, C22, C10, C13, C20]}} Each extra vessel coloured by fluoroscopy adds about the same, from about 22\% at 0 to about 88\% at 2 vessels, flattening out near the top because the predicted risk cannot go much higher; a flat rather than upsloping ST segment raises it to about 67\%, higher cholesterol raises it steadily from about 14\% at 193 mg/dL to about 37\% at 308 mg/dL, and chest pain brought on by exercise raises it to about 33\%.\,{\scriptsize\textcolor{NavyBlue}{[C01, C04, C07, C16]}} Through this sum alone, higher resting blood pressure raises it steadily, from about 14\% at 110 mm Hg to about 35\% at 155 mm Hg, older age raises it steadily, from about 19\% at 42 to about 24\% at 66 years, and a reversible defect on the thallium scan raises it to about 35\%; all three also act through the parts below.\,{\scriptsize\textcolor{NavyBlue}{[C08, C24, C15]}} Some rarer results also count (the model has little data on this, n = 2 to 17): through this sum alone a thallium result coded 4 raises the risk to about 83\% and a fixed defect on the thallium scan lowers it to about 11\%, and a downsloping ST segment raises it to about 27\%.\,{\scriptsize\textcolor{NavyBlue}{[C02, C18, C25]}} This part reads ST depression on exercise, the thallium scan and the resting ECG, and acts on 116 of 242 training patients, 77 of whom had heart disease; for the typical patient it acts on, the predicted risk is about 89\% without it and about 96\% with it.\,{\scriptsize\textcolor{NavyBlue}{[C40]}} Through it, more ST depression adds to the predicted risk only once it rises above about 0.9 mm, rising in an uneven way: for a typical patient from about 22\% at 0 mm to about 48\% at 2.8 mm.\,{\scriptsize\textcolor{NavyBlue}{[C33, C39]}} ST depression counts more with a reversible defect on the thallium scan (from about 36\% to about 80\%, against about 22\% to about 48\% otherwise) and with a normal resting ECG (from about 22\% to about 48\%, against about 22\% to about 34\% otherwise), in an uneven way.\,{\scriptsize\textcolor{NavyBlue}{[C36, C37]}} A reversible defect on the thallium scan also raises the risk more with a normal resting ECG, from about 22\% to about 47\% against about 22\% to about 36\% otherwise.\,{\scriptsize\textcolor{NavyBlue}{[C38]}} This part reads that signal together with resting blood pressure and chest pain type, and barely changes the predictions: for the typical patient it acts on, atypical angina raises the predicted risk through it from about 69\% to about 70\%, and non-anginal chest pain leaves it at about 69\%.\,{\scriptsize\textcolor{NavyBlue}{[C42, C44]}} This part acts on 85 of 242 training patients, 37 of whom had heart disease; for the typical patient it acts on, the predicted risk is about 28\% without it and about 30\% with it.\,{\scriptsize\textcolor{NavyBlue}{[C51]}} Through it, age adds nothing until about 53 years (between 53 and 66 years for different patients), then raises the risk steadily, for a typical patient from about 21\% at 42 to about 24\% at 66 years.\,{\scriptsize\textcolor{NavyBlue}{[C49, C50]}}
\paragraph{The whole model.} The model adds up its inputs into one score: the vessels coloured (each extra vessel adding about the same), the ST segment, sex and chest pain type each move the predicted risk on their own, and a chain of three parts adds to that sum: a signal from ST depression and the thallium scan, a blood pressure step that reads it, and an extra age term after that.\,{\scriptsize\textcolor{NavyBlue}{[C32, C01, C04, C12, C10]}} \textit{Evidence.} It is made up of the inputs that feed the score directly and three parts, the ST depression and thallium signal, the blood pressure step and the extra age term (each explained below); the first part feeds the second and the second feeds the third, and the score is their sum, which was checked on all 242 training patients.\,{\scriptsize\textcolor{NavyBlue}{[C32]}} ST depression on exercise and the resting ECG do not feed the score directly, so they act only through these parts.\,{\scriptsize\textcolor{NavyBlue}{[C32, C33]}} Changing one input at a time across the whole model, more ST depression raises a typical patient's predicted risk in an uneven way, from about 22\% at 0 mm to about 48\% at 2.8 mm, and a reversible defect on the thallium scan takes it to about 47\%, against about 35\% through the sum alone.\,{\scriptsize\textcolor{NavyBlue}{[C05, C06, C15]}} Across the whole model, higher resting blood pressure raises it in an uneven way, from about 14\% at 110 mm Hg to about 35\% at 155 mm Hg, and older age raises it in an uneven way, from about 18\% at 42 to about 26\% at 66 years.\,{\scriptsize\textcolor{NavyBlue}{[C09, C23]}} Across the whole model, a fixed defect on the thallium scan lowers it to about 11\% (the model has little data on this, n = 15), and a resting ECG other than normal barely changes it for a typical patient.\,{\scriptsize\textcolor{NavyBlue}{[C19, C26]}}
\paragraph{ST depression and thallium signal.} This part raises the predicted risk with ST depression on exercise once it rises above about 0.9 mm, in an uneven way, and more so with a reversible defect on the thallium scan or a normal resting ECG.\,{\scriptsize\textcolor{NavyBlue}{[C33, C39, C36, C37]}} \textit{Evidence.} It acts on 116 of 242 training patients, 77 of whom had heart disease; for the typical patient it acts on, the predicted risk is about 89\% without it and about 96\% with it.\,{\scriptsize\textcolor{NavyBlue}{[C40]}} Through it, more ST depression adds to the risk only once it rises above about 0.9 mm, rising in an uneven way: for a typical patient from about 22\% at 0 mm to about 48\% at 2.8 mm.\,{\scriptsize\textcolor{NavyBlue}{[C33, C39]}} ST depression counts more with a reversible defect on the thallium scan (from about 36\% to about 80\%, against about 22\% to about 48\% otherwise) and with a normal resting ECG (from about 22\% to about 48\%, against about 22\% to about 34\% otherwise), in an uneven way.\,{\scriptsize\textcolor{NavyBlue}{[C36, C37]}} Through it, a reversible defect on the thallium scan raises a typical patient's risk from about 22\% to about 32\%, and more with a normal resting ECG (to about 47\%, against about 36\% otherwise).\,{\scriptsize\textcolor{NavyBlue}{[C34, C38]}} Its formula reproduces the part exactly on all 242 training patients.\,{\scriptsize\textcolor{NavyBlue}{[C41]}}
\paragraph{Blood pressure step.} This part reads the ST depression and thallium signal together with resting blood pressure and chest pain type, and barely changes the predictions for the typical patient it acts on.\,{\scriptsize\textcolor{NavyBlue}{[C42, C44]}} \textit{Evidence.} For the typical patient it acts on, atypical angina raises the predicted risk through it from about 69\% to about 70\%, and non-anginal chest pain leaves it at about 69\%.\,{\scriptsize\textcolor{NavyBlue}{[C42, C44]}} Its formula reproduces the part exactly on all 242 training patients.\,{\scriptsize\textcolor{NavyBlue}{[C48]}}
\paragraph{Extra age term.} This part adds a little to the predicted risk for older patients: nothing until about 53 years, then steadily.\,{\scriptsize\textcolor{NavyBlue}{[C49, C50]}} \textit{Evidence.} It acts on 85 of 242 training patients, 37 of whom had heart disease; for the typical patient it acts on, the predicted risk is about 28\% without it and about 30\% with it.\,{\scriptsize\textcolor{NavyBlue}{[C51]}} Through it, age adds nothing until about 53 years (between 53 and 66 years for different patients), then raises the risk steadily, for a typical patient from about 21\% at 42 to about 24\% at 66 years.\,{\scriptsize\textcolor{NavyBlue}{[C49, C50]}} Its formula reproduces the part exactly on all 242 training patients.\,{\scriptsize\textcolor{NavyBlue}{[C52]}}
\paragraph{Worth checking against what is known.} Worth checking: the model treats fasting blood sugar above 120 mg/dL as lowering the risk of heart disease (for a typical patient from about 22\% to about 11\%).\,{\scriptsize\textcolor{NavyBlue}{[C17]}} Worth checking: patients with typical angina get the lowest predicted risk of all chest pain types, below patients with no chest pain symptoms.\,{\scriptsize\textcolor{NavyBlue}{[C11, C14, C21]}} Worth checking: the model treats a fixed defect on the thallium scan as lowering the risk compared with a normal scan, from only 15 training patients.\,{\scriptsize\textcolor{NavyBlue}{[C19]}} Worth checking: the model gives a higher predicted risk to the 2 training patients whose thallium result is coded 4, a code the dataset documentation does not define.\,{\scriptsize\textcolor{NavyBlue}{[C03]}} Worth checking: ST depression on exercise and a reversible thallium defect raise the predicted risk more when the resting ECG is normal than when it is not.\,{\scriptsize\textcolor{NavyBlue}{[C37, C38]}}
\paragraph{Claims.}
\par{\tiny\noindent
\begin{tabularx}{\linewidth}{@{}l>{\raggedright\arraybackslash}X>{\raggedright\arraybackslash}p{0.11\linewidth}>{\raggedright\arraybackslash}p{0.16\linewidth}>{\raggedright\arraybackslash}p{0.05\linewidth}>{\raggedright\arraybackslash}p{0.15\linewidth}l>{\raggedright\arraybackslash}p{0.08\linewidth}@{}}
\toprule
Claim & What is checked & Risk, typical patient & Estimate (log-odds; OR); range across patients & n & Test check & Verified & Words \\
\midrule
C01 & direction: number of vessels coloured by fluoroscopy: 0 vessels to 2 vessels (whole model) \newline\textit{shape: steady (equal steps); 100\% of patients same; verified; risk saturates near 100\%} & 22\% to 88\% & log-odds 3.32, OR 27.6; range across patients 3.32 to 3.32 & 242 & train; test unclear (AUC gain 0.039 [-0.026, 0.111], n=61) & \textcolor{OliveGreen}{yes} & on the training data only \\
C02 & direction: thallium stress test result: a normal thallium scan to thallium stress test result coded 4 (meaning not recorded) (direct contribution only, not through any part) \newline\textit{shape: categorical; highest to lowest: thallium stress test result coded 4 (meaning not recorded), a reversible defect on the thallium scan, a normal thallium scan, a fixed defect on the thallium scan} & 22\% to 83\% & log-odds 2.85, OR 17.3; range across patients 2.85 to 2.85 & 135, 2 (small n) & train; test unclear (AUC gain 0.012 [-0.009, 0.035], n=61) & \textcolor{OliveGreen}{yes} & the model has little data on this; on the training data only \\
C03 & direction: thallium stress test result: a normal thallium scan to thallium stress test result coded 4 (meaning not recorded) (whole model) \newline\textit{shape: categorical; highest to lowest: thallium stress test result coded 4 (meaning not recorded), a reversible defect on the thallium scan, a normal thallium scan, a fixed defect on the thallium scan} & 22\% to 83\% & log-odds 2.85, OR 17.3; range across patients 2.85 to 2.85 & 135, 2 (small n) & train; test unclear (AUC gain 0.017 [-0.010, 0.052], n=61) & \textcolor{OliveGreen}{yes} & the model has little data on this; on the training data only \\
C04 & direction: slope of the ST segment at peak exercise: an upsloping ST segment to a flat ST segment (whole model) \newline\textit{shape: categorical; highest to lowest: a flat ST segment, a downsloping ST segment, an upsloping ST segment} & 22\% to 67\% & log-odds 2.02, OR 7.53; range across patients 2.02 to 2.02 & 115, 110 & train; test disagrees (AUC gain -0.019 [-0.070, 0.025], n=61) & \textcolor{OliveGreen}{yes} & on the training data only \\
C05 & direction: ST depression on exercise: 0.0 mm to 2.8 mm (whole model) \newline\textit{shape: accelerating; 48\% of patients same; fallback} & 22\% to 48\% & log-odds 1.16, OR 3.20; range across patients 0.55 to 1.91 & 242 & train; test unclear (AUC gain 0.008 [-0.011, 0.031], n=61) & \textcolor{OliveGreen}{yes} & on the training data only \\
C06 & direction: thallium stress test result: a normal thallium scan to a reversible defect on the thallium scan (whole model) \newline\textit{shape: categorical; highest to lowest: thallium stress test result coded 4 (meaning not recorded), a reversible defect on the thallium scan, a normal thallium scan, a fixed defect on the thallium scan} & 22\% to 47\% & log-odds 1.01, OR 2.73; range across patients 0.67 to 1.44 & 135, 90 & train; test unclear (AUC gain 0.017 [-0.010, 0.052], n=61) & \textcolor{OliveGreen}{yes} & on the training data only \\
C07 & direction: serum cholesterol: 193 mg/dL (5.0 mmol/L) to 308 mg/dL (8.0 mmol/L) (whole model) \newline\textit{shape: steady; 100\% of patients same; verified} & 14\% to 37\% & log-odds 1.30, OR 3.68; range across patients 1.30 to 1.30 & 242 & train; test disagrees (AUC gain -0.022 [-0.063, 0.008], n=61) & \textcolor{OliveGreen}{yes} & on the training data only \\
C08 & direction: resting blood pressure: 110 mm Hg to 155 mm Hg (direct contribution only, not through any part) \newline\textit{shape: steady; 100\% of patients same; verified} & 14\% to 35\% & log-odds 1.19, OR 3.28; range across patients 1.19 to 1.19 & 242 & train; test unclear (AUC gain 0.000 [-0.024, 0.022], n=61) & \textcolor{OliveGreen}{yes} & on the training data only \\
C09 & direction: resting blood pressure: 110 mm Hg to 155 mm Hg (whole model) \newline\textit{shape: steady; 43\% of patients same; fallback} & 14\% to 35\% & log-odds 1.04, OR 2.82; range across patients 0.84 to 1.21 & 242 & train; test unclear (AUC gain 0.001 [-0.022, 0.023], n=61) & \textcolor{OliveGreen}{yes} & on the training data only \\
C10 & direction: chest pain type: no chest pain (asymptomatic) to typical angina (direct contribution only, not through any part) \newline\textit{shape: categorical; highest to lowest: no chest pain (asymptomatic), atypical angina, non-anginal chest pain, typical angina} & 22\% to 2\% & log-odds -2.48, OR 0.08; range across patients -2.48 to -2.48 & 108, 22 & train; test agrees (AUC gain 0.040 [0.008, 0.079], n=61) & \textcolor{OliveGreen}{yes} & clearly \\
C11 & direction: chest pain type: no chest pain (asymptomatic) to typical angina (whole model) \newline\textit{shape: categorical; highest to lowest: no chest pain (asymptomatic), atypical angina, non-anginal chest pain, typical angina} & 22\% to 2\% & log-odds -2.48, OR 0.08; range across patients -2.48 to -2.48 & 108, 22 & train; test agrees (AUC gain 0.040 [0.008, 0.079], n=61) & \textcolor{OliveGreen}{yes} & clearly \\
C12 & direction: sex: male to female (whole model) \newline\textit{shape: categorical; highest to lowest: male, female} & 22\% to 3\% & log-odds -2.26, OR 0.10; range across patients -2.26 to -2.26 & 159, 83 & train; test unclear (AUC gain 0.024 [-0.008, 0.064], n=61) & \textcolor{OliveGreen}{yes} & on the training data only \\
\bottomrule
\end{tabularx}}
\par{\tiny\noindent
\begin{tabularx}{\linewidth}{@{}l>{\raggedright\arraybackslash}X>{\raggedright\arraybackslash}p{0.11\linewidth}>{\raggedright\arraybackslash}p{0.16\linewidth}>{\raggedright\arraybackslash}p{0.05\linewidth}>{\raggedright\arraybackslash}p{0.15\linewidth}l>{\raggedright\arraybackslash}p{0.08\linewidth}@{}}
\toprule
Claim & What is checked & Risk, typical patient & Estimate (log-odds; OR); range across patients & n & Test check & Verified & Words \\
\midrule
C13 & direction: chest pain type: no chest pain (asymptomatic) to non-anginal chest pain (direct contribution only, not through any part) \newline\textit{shape: categorical; highest to lowest: no chest pain (asymptomatic), atypical angina, non-anginal chest pain, typical angina} & 22\% to 5\% & log-odds -1.73, OR 0.18; range across patients -1.73 to -1.73 & 108, 70 & train; test agrees (AUC gain 0.040 [0.008, 0.079], n=61) & \textcolor{OliveGreen}{yes} & clearly \\
C14 & direction: chest pain type: no chest pain (asymptomatic) to non-anginal chest pain (whole model) \newline\textit{shape: categorical; highest to lowest: no chest pain (asymptomatic), atypical angina, non-anginal chest pain, typical angina} & 22\% to 5\% & log-odds -1.73, OR 0.18; range across patients -1.73 to -1.72 & 108, 70 & train; test agrees (AUC gain 0.040 [0.008, 0.079], n=61) & \textcolor{OliveGreen}{yes} & clearly \\
C15 & direction: thallium stress test result: a normal thallium scan to a reversible defect on the thallium scan (direct contribution only, not through any part) \newline\textit{shape: categorical; highest to lowest: thallium stress test result coded 4 (meaning not recorded), a reversible defect on the thallium scan, a normal thallium scan, a fixed defect on the thallium scan} & 22\% to 35\% & log-odds 0.67, OR 1.95; range across patients 0.67 to 0.67 & 135, 90 & train; test unclear (AUC gain 0.012 [-0.009, 0.035], n=61) & \textcolor{OliveGreen}{yes} & on the training data only \\
C16 & direction: exercise-induced angina: no chest pain on exercise to chest pain brought on by exercise (whole model) \newline\textit{shape: categorical; highest to lowest: chest pain brought on by exercise, no chest pain on exercise} & 22\% to 33\% & log-odds 0.60, OR 1.82; range across patients 0.60 to 0.60 & 160, 82 & train; test unclear (AUC gain 0.012 [-0.003, 0.032], n=61) & \textcolor{OliveGreen}{yes} & on the training data only \\
C17 & direction: fasting blood sugar: fasting blood sugar of 120 mg/dL or less to fasting blood sugar above 120 mg/dL (whole model) \newline\textit{shape: categorical; highest to lowest: fasting blood sugar of 120 mg/dL or less, fasting blood sugar above 120 mg/dL} & 22\% to 11\% & log-odds -0.81, OR 0.44; range across patients -0.81 to -0.81 & 209, 33 & train; test disagrees (AUC gain -0.002 [-0.020, 0.015], n=61) & \textcolor{OliveGreen}{yes} & on the training data only \\
C18 & direction: thallium stress test result: a normal thallium scan to a fixed defect on the thallium scan (direct contribution only, not through any part) \newline\textit{shape: categorical; highest to lowest: thallium stress test result coded 4 (meaning not recorded), a reversible defect on the thallium scan, a normal thallium scan, a fixed defect on the thallium scan} & 22\% to 11\% & log-odds -0.77, OR 0.46; range across patients -0.77 to -0.77 & 135, 15 (small n) & train; test unclear (AUC gain 0.012 [-0.009, 0.035], n=61) & \textcolor{OliveGreen}{yes} & the model has little data on this; on the training data only \\
C19 & direction: thallium stress test result: a normal thallium scan to a fixed defect on the thallium scan (whole model) \newline\textit{shape: categorical; highest to lowest: thallium stress test result coded 4 (meaning not recorded), a reversible defect on the thallium scan, a normal thallium scan, a fixed defect on the thallium scan} & 22\% to 11\% & log-odds -0.77, OR 0.46; range across patients -0.77 to -0.77 & 135, 15 (small n) & train; test unclear (AUC gain 0.017 [-0.010, 0.052], n=61) & \textcolor{OliveGreen}{yes} & the model has little data on this; on the training data only \\
C20 & direction: chest pain type: no chest pain (asymptomatic) to atypical angina (direct contribution only, not through any part) \newline\textit{shape: categorical; highest to lowest: no chest pain (asymptomatic), atypical angina, non-anginal chest pain, typical angina} & 22\% to 11\% & log-odds -0.76, OR 0.47; range across patients -0.76 to -0.76 & 108, 42 & train; test agrees (AUC gain 0.040 [0.008, 0.079], n=61) & \textcolor{OliveGreen}{yes} & clearly \\
C21 & direction: chest pain type: no chest pain (asymptomatic) to atypical angina (whole model) \newline\textit{shape: categorical; highest to lowest: no chest pain (asymptomatic), atypical angina, non-anginal chest pain, typical angina} & 22\% to 11\% & log-odds -0.72, OR 0.49; range across patients -0.76 to -0.68 & 108, 42 & train; test agrees (AUC gain 0.040 [0.008, 0.079], n=61) & \textcolor{OliveGreen}{yes} & clearly \\
C22 & direction: maximum heart rate: 118 beats per minute to 175 beats per minute (whole model) \newline\textit{shape: steady; 100\% of patients same; verified} & 27\% to 18\% & log-odds -0.53, OR 0.59; range across patients -0.53 to -0.53 & 242 & train; test unclear (AUC gain 0.006 [0.000, 0.022], n=61) & \textcolor{OliveGreen}{yes} & on the training data only \\
C23 & direction: age: 42 years to 66 years (whole model) \newline\textit{shape: accelerating; 67\% of patients same; fallback} & 18\% to 26\% & log-odds 0.41, OR 1.50; range across patients 0.31 to 0.49 & 242 & train; test disagrees (AUC gain -0.003 [-0.017, 0.013], n=61) & \textcolor{OliveGreen}{yes} & on the training data only \\
C24 & direction: age: 42 years to 66 years (direct contribution only, not through any part) \newline\textit{shape: steady; 100\% of patients same; verified} & 19\% to 24\% & log-odds 0.31, OR 1.36; range across patients 0.31 to 0.31 & 242 & train; test disagrees (AUC gain -0.004 [-0.018, 0.012], n=61) & \textcolor{OliveGreen}{yes} & on the training data only \\
\bottomrule
\end{tabularx}}
\par{\tiny\noindent
\begin{tabularx}{\linewidth}{@{}l>{\raggedright\arraybackslash}X>{\raggedright\arraybackslash}p{0.11\linewidth}>{\raggedright\arraybackslash}p{0.16\linewidth}>{\raggedright\arraybackslash}p{0.05\linewidth}>{\raggedright\arraybackslash}p{0.15\linewidth}l>{\raggedright\arraybackslash}p{0.08\linewidth}@{}}
\toprule
Claim & What is checked & Risk, typical patient & Estimate (log-odds; OR); range across patients & n & Test check & Verified & Words \\
\midrule
C25 & direction: slope of the ST segment at peak exercise: an upsloping ST segment to a downsloping ST segment (whole model) \newline\textit{shape: categorical; highest to lowest: a flat ST segment, a downsloping ST segment, an upsloping ST segment} & 22\% to 27\% & log-odds 0.28, OR 1.32; range across patients 0.28 to 0.28 & 115, 17 (small n) & train; test disagrees (AUC gain -0.019 [-0.070, 0.025], n=61) & \textcolor{OliveGreen}{yes} & the model has little data on this; on the training data only \\
C26 & direction: resting ECG: a normal resting ECG to resting ECG other than a normal resting ECG (whole model) \newline\textit{shape: categorical; highest to lowest: a normal resting ECG, resting ECG other than a normal resting ECG} & 22\% to 22\% & log-odds -0.23, OR 0.80; range across patients -0.57 to 0.00 & 121, 121 & train; test disagrees (AUC gain -0.002 [-0.015, 0.009], n=61) & \textcolor{OliveGreen}{yes} & slightly; on the training data only \\
C27 & magnitude: number of vessels coloured by fluoroscopy: 0 vessels to 2 vessels, typical patient \newline\textit{shape: steady (equal steps); 100\% of patients same; verified; risk saturates near 100\%} & 22\% to 88\% & log-odds 3.32, OR 27.6 & 242 & train; test unclear (AUC gain 0.039 [-0.026, 0.111], n=61) & \textcolor{OliveGreen}{yes} & on the training data only \\
C28 & magnitude: slope of the ST segment at peak exercise: an upsloping ST segment to a flat ST segment, typical patient \newline\textit{shape: categorical; highest to lowest: a flat ST segment, a downsloping ST segment, an upsloping ST segment} & 22\% to 67\% & log-odds 2.02, OR 7.53 & 115, 110 & train; test disagrees (AUC gain -0.019 [-0.070, 0.025], n=61) & \textcolor{OliveGreen}{yes} & on the training data only \\
C29 & magnitude: ST depression on exercise: 0.0 mm to 2.8 mm, typical patient \newline\textit{shape: accelerating; 48\% of patients same; fallback} & 22\% to 48\% & log-odds 1.21, OR 3.36 & 242 & train; test unclear (AUC gain 0.008 [-0.011, 0.031], n=61) & \textcolor{OliveGreen}{yes} & on the training data only \\
C30 & magnitude: thallium stress test result: a normal thallium scan to a reversible defect on the thallium scan, typical patient \newline\textit{shape: categorical; highest to lowest: thallium stress test result coded 4 (meaning not recorded), a reversible defect on the thallium scan, a normal thallium scan, a fixed defect on the thallium scan} & 22\% to 47\% & log-odds 1.19, OR 3.28 & 135, 90 & train; test unclear (AUC gain 0.017 [-0.010, 0.052], n=61) & \textcolor{OliveGreen}{yes} & on the training data only \\
C31 & magnitude: serum cholesterol: 193 mg/dL (5.0 mmol/L) to 308 mg/dL (8.0 mmol/L), typical patient \newline\textit{shape: steady; 100\% of patients same; verified} & 14\% to 37\% & log-odds 1.30, OR 3.68 & 242 & train; test disagrees (AUC gain -0.022 [-0.063, 0.008], n=61) & \textcolor{OliveGreen}{yes} & on the training data only \\
C32 & structure: made of ischaemia\_\allowbreak{}gate, bp\_\allowbreak{}offset\_\allowbreak{}gate, age\_\allowbreak{}gate\_\allowbreak{}without\_\allowbreak{}ischaemia and inputs age, sex, cp, trestbps, chol, fbs, thalach, exang, slope, ca, thal; added; ann\_\allowbreak{}32 feeds ann\_\allowbreak{}33; ann\_\allowbreak{}33 feeds ann\_\allowbreak{}34 &  & log-odds n/a, OR n/a & 242 & train & \textcolor{OliveGreen}{yes} &  \\
C33 & direction: ST depression on exercise: 0.0 mm to 2.8 mm (through this part) \newline\textit{shape: accelerating; 48\% of patients same; fallback} & 22\% to 48\% & log-odds 1.16, OR 3.20; range across patients 0.55 to 1.91 & 242 & train; test unclear (AUC gain 0.008 [-0.011, 0.031], n=61) & \textcolor{OliveGreen}{yes} & on the training data only \\
C34 & direction: thallium stress test result: thallium stress test result other than a reversible defect on the thallium scan to a reversible defect on the thallium scan (through this part) \newline\textit{shape: categorical; highest to lowest: a reversible defect on the thallium scan, thallium stress test result other than a reversible defect on the thallium scan} & 22\% to 32\% & log-odds 0.34, OR 1.40; range across patients 0.00 to 0.77 & 152, 90 & train; test unclear (AUC gain 0.005 [-0.009, 0.025], n=61) & \textcolor{OliveGreen}{yes} & on the training data only \\
C35 & direction: resting ECG: a normal resting ECG to resting ECG other than a normal resting ECG (through this part) \newline\textit{shape: categorical; highest to lowest: a normal resting ECG, resting ECG other than a normal resting ECG} & 96\% to 94\% (typical patient this part acts on) & log-odds -0.23, OR 0.80; range across patients -0.57 to 0.00 & 121, 121 & train; test disagrees (AUC gain -0.002 [-0.015, 0.009], n=61) & \textcolor{OliveGreen}{yes} & slightly; on the training data only \\
C36 & interaction: effect of ST depression on exercise (0.0 mm to 2.8 mm) by thallium stress test result: a reversible defect on the thallium scan minus thallium stress test result other than a reversible defect on the thallium scan; stronger in a reversible defect on the thallium scan \newline\textit{shape: steady; 37\% of patients same; fallback} & a reversible defect on the thallium scan: 36\% to 80\%; thallium stress test result other than a reversible defect on the thallium scan: 22\% to 48\% & log-odds difference 0.68, OR 1.98 & 90, 152 & train; test agrees (AUC gain 0.005 [-0.009, 0.025], n=61); group difference on test 0.73 [0.59, 0.86] & \textcolor{OliveGreen}{yes} & clearly \\
\bottomrule
\end{tabularx}}
\par{\tiny\noindent
\begin{tabularx}{\linewidth}{@{}l>{\raggedright\arraybackslash}X>{\raggedright\arraybackslash}p{0.11\linewidth}>{\raggedright\arraybackslash}p{0.16\linewidth}>{\raggedright\arraybackslash}p{0.05\linewidth}>{\raggedright\arraybackslash}p{0.15\linewidth}l>{\raggedright\arraybackslash}p{0.08\linewidth}@{}}
\toprule
Claim & What is checked & Risk, typical patient & Estimate (log-odds; OR); range across patients & n & Test check & Verified & Words \\
\midrule
C37 & interaction: effect of ST depression on exercise (0.0 mm to 2.8 mm) by resting ECG: resting ECG other than a normal resting ECG minus a normal resting ECG; stronger in a normal resting ECG \newline\textit{shape: accelerating; 33\% of patients same; fallback} & resting ECG other than a normal resting ECG: 22\% to 34\%; a normal resting ECG: 22\% to 48\% & log-odds difference -0.59, OR 0.55 & 121, 121 & train; test agrees (AUC gain -0.002 [-0.015, 0.009], n=61); group difference on test -0.59 [-0.76, -0.40] & \textcolor{OliveGreen}{yes} & clearly \\
C38 & interaction: effect of thallium stress test result (thallium stress test result other than a reversible defect on the thallium scan to a reversible defect on the thallium scan) by resting ECG: resting ECG other than a normal resting ECG minus a normal resting ECG; stronger in a normal resting ECG \newline\textit{shape: categorical; highest to lowest: a reversible defect on the thallium scan, thallium stress test result other than a reversible defect on the thallium scan} & resting ECG other than a normal resting ECG: 22\% to 36\%; a normal resting ECG: 22\% to 47\% & log-odds difference -0.11, OR 0.90 & 121, 121 & train; test agrees (AUC gain -0.002 [-0.015, 0.009], n=61); group difference on test -0.18 [-0.34, -0.03] & \textcolor{OliveGreen}{yes} & clearly \\
C39 & threshold: ST depression on exercise: fold at about 0.9 mm (starts) \newline\textit{shape: accelerating; 55\% of patients same; fallback} & 22\% to 46\% (typical patient with a fold) & fold 0.93 mm (IQR 0.93 to 1.72); past it log-odds 1.50, OR 4.47; range across patients 1.00 to 1.86 & 198, 72 & train; test unclear (AUC gain 0.008 [-0.011, 0.031], n=61) & \textcolor{OliveGreen}{yes} & on the training data only \\
C40 & subgroup: who this part acts on: 116 of 242 patients & 89\% to 96\% (typical patient this part acts on (risk without the part, then with it)) & share 0.48; log-odds 0.77, OR 2.17; range across patients 0.08 to 1.75 & 116, 126 & train; test unclear (AUC gain 0.016 [-0.005, 0.048], n=61) & \textcolor{OliveGreen}{yes} & on the training data only \\
C41 & structure: its formula reproduces it on 242 training patients (largest difference 1.77636e-15) &  & log-odds 2e-15, OR n/a & 242 & train & \textcolor{OliveGreen}{yes} &  \\
C42 & direction: chest pain type: chest pain type other than atypical angina, non-anginal chest pain to atypical angina (through this part) \newline\textit{shape: categorical; highest to lowest: atypical angina, non-anginal chest pain, chest pain type other than atypical angina, non-anginal chest pain} & 69\% to 70\% (typical patient this part acts on) & log-odds 0.04, OR 1.04; range across patients 0.00 to 0.08 & 130, 42 & train; test unclear (AUC gain 0.000 [0.000, 0.000], n=61) & \textcolor{OliveGreen}{yes} & slightly; on the training data only \\
C43 & direction: resting blood pressure: 110 mm Hg to 155 mm Hg (through this part) \newline\textit{shape: hinge at 114 mm Hg; 45\% of patients same; fallback} & 22\% to 22\% & log-odds -0.15, OR 0.86; range across patients -0.34 to 0.02 & 242 & train; test disagrees (AUC gain -0.003 [-0.015, 0.006], n=61) & \textcolor{BrickRed}{no} & slightly; on the training data only \\
C44 & direction: chest pain type: chest pain type other than atypical angina, non-anginal chest pain to non-anginal chest pain (through this part) \newline\textit{shape: categorical; highest to lowest: atypical angina, non-anginal chest pain, chest pain type other than atypical angina, non-anginal chest pain} & 69\% to 69\% (typical patient this part acts on) & log-odds 0.003, OR 1.00; range across patients 0.00 to 0.006 & 130, 70 & train; test unclear (AUC gain 0.000 [0.000, 0.000], n=61) & \textcolor{OliveGreen}{yes} & slightly; on the training data only \\
C45 & interaction: effect of resting blood pressure (110 mm Hg to 155 mm Hg) by chest pain type: atypical angina minus chest pain type other than atypical angina, non-anginal chest pain; stronger in atypical angina \newline\textit{shape: accelerating; 90\% of patients same; verified} & atypical angina: 8\% to 20\%; chest pain type other than atypical angina, non-anginal chest pain: 14\% to 35\% & log-odds difference 0.04, OR 1.04 & 42, 130 & train; test unclear (AUC gain 0.000 [0.000, 0.000], n=61); group difference on test 0.06 [-0.03, 0.14] & \textcolor{BrickRed}{no} & on the training data only \\
C46 & threshold: resting blood pressure: fold at about 128 mm Hg (stops) \newline\textit{shape: several pieces; 14\% of patients same; fallback} & 22\% to 21\% (typical patient with a fold) & fold 128 mm Hg (IQR 120 to 138); past it log-odds -0.16, OR 0.85; range across patients -0.42 to 0.03 & 203, 101 & train; test disagrees (AUC gain -0.003 [-0.015, 0.006], n=61) & \textcolor{BrickRed}{no} & slightly; on the training data only \\
C47 & subgroup: who this part acts on: 140 of 242 patients & 65\% to 69\% (typical patient this part acts on (risk without the part, then with it)) & share 0.58; log-odds 0.33, OR 1.39; range across patients -0.01 to 0.84 & 140, 102 & train; test unclear (AUC gain 0.006 [-0.004, 0.022], n=61) & \textcolor{BrickRed}{no} & slightly; on the training data only \\
C48 & structure: its formula reproduces it on 242 training patients (largest difference 8.88178e-16) &  & log-odds 9e-16, OR n/a & 242 & train & \textcolor{OliveGreen}{yes} &  \\
\bottomrule
\end{tabularx}}
\par{\tiny\noindent
\begin{tabularx}{\linewidth}{@{}l>{\raggedright\arraybackslash}X>{\raggedright\arraybackslash}p{0.11\linewidth}>{\raggedright\arraybackslash}p{0.16\linewidth}>{\raggedright\arraybackslash}p{0.05\linewidth}>{\raggedright\arraybackslash}p{0.15\linewidth}l>{\raggedright\arraybackslash}p{0.08\linewidth}@{}}
\toprule
Claim & What is checked & Risk, typical patient & Estimate (log-odds; OR); range across patients & n & Test check & Verified & Words \\
\midrule
C49 & direction: age: 42 years to 66 years (through this part) \newline\textit{shape: hinge at 53 years; 99\% of patients same; verified; one point with C50 (median fold 53 years, typical patient 53 years, 10th-90th percentile 53 years to 66 years, wide)} & 21\% to 24\% & log-odds 0.10, OR 1.11; range across patients 0.00 to 0.18 & 242 & train; test unclear (AUC gain 0.000 [-0.01, 0.011], n=61) & \textcolor{OliveGreen}{yes} & slightly; on the training data only \\
C50 & threshold: age: fold at about 53 years (starts) \newline\textit{shape: hinge at 53 years; 99\% of patients same; verified; one point with C50 (median fold 53 years, typical patient 53 years, 10th-90th percentile 53 years to 66 years, wide)} & 20\% to 24\% (typical patient with a fold) & fold 52.7 years (IQR 52.7 to 58.6); past it log-odds 0.17, OR 1.18; range across patients 0.03 to 0.21 & 174, 85 & train; test unclear (AUC gain 0.000 [-0.01, 0.011], n=61) & \textcolor{OliveGreen}{yes} & slightly; on the training data only \\
C51 & subgroup: who this part acts on: 85 of 242 patients & 28\% to 30\% (typical patient this part acts on (risk without the part, then with it)) & share 0.35; log-odds 0.10, OR 1.11; range across patients 0.01 to 0.20 & 85, 157 & train; test disagrees (AUC gain -0.002 [-0.01, 0.004], n=61) & \textcolor{OliveGreen}{yes} & slightly; on the training data only \\
C52 & structure: its formula reproduces it on 242 training patients (largest difference 1.11022e-16) &  & log-odds 1e-16, OR n/a & 242 & train & \textcolor{OliveGreen}{yes} &  \\
\bottomrule
\end{tabularx}}
\paragraph{Not supported by the model.}
\begin{itemize}\setlength\itemsep{0pt}\small
\item C43: direction claim about resting blood pressure in part ann\_\allowbreak{}33: not supported by the model (17\% of patients have the opposite sign (allowed: 10\%)).
\item C45: interaction claim about resting blood pressure and chest pain type in part ann\_\allowbreak{}33: not supported by the model (the group difference is below 0.1 log-odds; the part's own counterfactual difference has the opposite sign).
\item C46: threshold claim about resting blood pressure in part ann\_\allowbreak{}33: not supported by the model (the effect past the fold is too small or not consistent across patients).
\item C47: subgroup claim about chest pain type and resting blood pressure in part ann\_\allowbreak{}33: not supported by the model (its effect goes both ways across the patients it acts on).
\end{itemize}

\subsection{Coverage}\label{app:heart-s5-c1:coverage}
Coverage is computed by the current code (the function behind the \texttt{get\_coverage} tool) under these conventions. \emph{Candidate nodes} are all nodes of the explained model state except the output nodes. A candidate node is \emph{covered} when some annotation of the hierarchy (a leaf, or a composition together with its descendants) contains it and all its enabled outgoing edges. \emph{Exact Null} parts, which have no effect on any output for any input, count as covered: nodes with no path to an output over enabled, non-zero-weight edges (disconnected inputs and dead ends), and edges that have zero weight or touch such a node. Compositional coverage is the fraction of required compositions present: every annotation except a single global root containing every output must have a parent; a parentless, childless annotation made only of Null nodes is exempt.

\begin{itemize}
\item \textbf{Null, disconnected input (4):} \texttt{-11} (fbs:1), \texttt{-13} (restecg:1), \texttt{-14} (restecg:2), \texttt{-3} (sex:1).
\item \textbf{Uncovered nodes:} none.
\item \textbf{Uncovered edges:} none.
\item \textbf{Missing compositions:} none.
\end{itemize}

\subsection{Annotation hierarchy}\label{app:heart-s5-c1:hierarchy}
The hierarchy has 1 root; the global root is heart\_\allowbreak{}disease\_\allowbreak{}logit. Tags (A1, A2, \ldots) are used in the figures.
\begingroup\parindent=0pt\parskip=1pt
\hspace*{0.0em}$\triangleright$~A1: heart\_\allowbreak{}disease\_\allowbreak{}logit (composition, \cref{app:heart-s5-c1:ann-heart-disease-logit})\par
\hspace*{1.5em}$\triangleright$~A2: ischaemia\_\allowbreak{}gate (leaf, \cref{app:heart-s5-c1:ann-ischaemia-gate})\par
\hspace*{1.5em}$\triangleright$~A3: bp\_\allowbreak{}offset\_\allowbreak{}gate (leaf, \cref{app:heart-s5-c1:ann-bp-offset-gate})\par
\hspace*{1.5em}$\triangleright$~A4: age\_\allowbreak{}gate\_\allowbreak{}without\_\allowbreak{}ischaemia (leaf, \cref{app:heart-s5-c1:ann-age-gate-without-ischaemia})\par
\endgroup

\subsection{Network figures}\label{app:heart-s5-c1:figures}
\begin{figure}[htbp]
\centering
\input{\explaneatcertdir{heart-s5-c1}heart-s5-c1-fig-full.tex}
\caption{heart-s5-c1: full view of the explained model state.}
\label{app:heart-s5-c1:fig-full}
\end{figure}
\begin{figure}[htbp]
\centering
\input{\explaneatcertdir{heart-s5-c1}heart-s5-c1-fig-level-0.tex}
\caption{heart-s5-c1: collapsed view at hierarchy level 0 (fully recomposed).}
\label{app:heart-s5-c1:fig-level-0}
\end{figure}
\begin{figure}[htbp]
\centering
\input{\explaneatcertdir{heart-s5-c1}heart-s5-c1-fig-level-1.tex}
\caption{heart-s5-c1: collapsed view at hierarchy level 1.}
\label{app:heart-s5-c1:fig-level-1}
\end{figure}

\subsection{Refactoring operations}\label{app:heart-s5-c1:operations}
The explained model state is the phenotype after these operations, replayed in order. All are function-preserving; the largest change in the model output over the 242 training rows is $2.22\times 10^{-16}$ (tolerance $1.00\times 10^{-9}$).
\begin{itemize}
\item \#26: split \texttt{-26} into \texttt{-26\_\allowbreak{}a}, \texttt{-26\_\allowbreak{}b}.
\item \#27: split \texttt{-6} into \texttt{-6\_\allowbreak{}a}, \texttt{-6\_\allowbreak{}b}.
\item \#28: split \texttt{-8} into \texttt{-8\_\allowbreak{}a}, \texttt{-8\_\allowbreak{}b}.
\item \#29: split \texttt{-1} into \texttt{-1\_\allowbreak{}a}, \texttt{-1\_\allowbreak{}b}.
\item \#30: identity node \texttt{identity\_\allowbreak{}290} before \texttt{290} (1 connection(s)).
\item \#31: identity node \texttt{identity\_\allowbreak{}1987} before \texttt{1987} (1 connection(s)).
\end{itemize}

\subsection{Annotations}\label{app:heart-s5-c1:annotations}
\subsubsection{\texorpdfstring{A1: heart\_\allowbreak{}disease\_\allowbreak{}logit}{A1: heart disease logit}}\label{app:heart-s5-c1:ann-heart-disease-logit}
% annotation id ann_35, operation seq 35
\textbf{Kind:} Composition annotation (depth 0); children A2 (\cref{app:heart-s5-c1:ann-ischaemia-gate}), A3 (\cref{app:heart-s5-c1:ann-bp-offset-gate}), A4 (\cref{app:heart-s5-c1:ann-age-gate-without-ischaemia}).

\begin{itemize}
\item \textbf{Entry nodes:} \texttt{-1\_\allowbreak{}a} (age\_\allowbreak{}a), \texttt{-2} (sex:0), \texttt{-3} (sex:1), \texttt{-4} (cp:1), \texttt{-6\_\allowbreak{}a} (cp:3\_\allowbreak{}a), \texttt{-7} (cp:4), \texttt{-8\_\allowbreak{}a} (trestbps\_\allowbreak{}a), \texttt{-9} (chol), \texttt{-10} (fbs:0), \texttt{-11} (fbs:1), \texttt{-13} (restecg:1), \texttt{-14} (restecg:2), \texttt{-15} (thalach), \texttt{-16} (exang:0), \texttt{-17} (exang:1), \texttt{-19} (slope:1), \texttt{-20} (slope:2), \texttt{-21} (slope:3), \texttt{-22} (ca), \texttt{-23} (thal:3), \texttt{-24} (thal:4), \texttt{-25} (thal:6), \texttt{-26\_\allowbreak{}a} (thal:7\_\allowbreak{}a).
\item \textbf{Exit nodes:} \texttt{0}.
\item \textbf{Internal nodes:} none.
\item \textbf{Edges (weight):} \texttt{-2}$\to$\texttt{0}~(-1.07), \texttt{-4}$\to$\texttt{0}~(-0.493), \texttt{-7}$\to$\texttt{0}~(0.377), \texttt{-9}$\to$\texttt{0}~(0.546), \texttt{-10}$\to$\texttt{0}~(0.278), \texttt{-15}$\to$\texttt{0}~(-0.207), \texttt{-16}$\to$\texttt{0}~(0.149), \texttt{-17}$\to$\texttt{0}~(0.431), \texttt{-19}$\to$\texttt{0}~(-0.855), \texttt{-20}$\to$\texttt{0}~(0.153), \texttt{-21}$\to$\texttt{0}~(-0.366), \texttt{-22}$\to$\texttt{0}~(1.48), \texttt{-23}$\to$\texttt{0}~(-0.41), \texttt{-24}$\to$\texttt{0}~(0.183), \texttt{-25}$\to$\texttt{0}~(-0.384), \texttt{204}$\to$\texttt{0}~(0.143), \texttt{290}$\to$\texttt{0}~(0.17), \texttt{1987}$\to$\texttt{0}~(0.143), \texttt{-26\_\allowbreak{}a}$\to$\texttt{0}~(-0.0766), \texttt{-6\_\allowbreak{}a}$\to$\texttt{0}~(-0.44), \texttt{-8\_\allowbreak{}a}$\to$\texttt{0}~(0.48), \texttt{-1\_\allowbreak{}a}$\to$\texttt{0}~(0.117), \texttt{204}$\to$\texttt{identity\_\allowbreak{}290}~(1.08), \texttt{290}$\to$\texttt{identity\_\allowbreak{}1987}~(-1.59).
\item \textbf{With its descendants:} 36 nodes.
\item \textbf{Covers} (coverage, own and as part of the hierarchy): 35 candidate nodes.
\end{itemize}

\textbf{Hypothesis.} Root: p = sigmoid(z), z = -0.504 + LINEAR CORE + 0.143*ischaemia\_\allowbreak{}gate + 0.170*bp\_\allowbreak{}offset\_\allowbreak{}gate + 0.143*\allowbreak{}age\_\allowbreak{}gate\_\allowbreak{}without\_\allowbreak{}ischaemia. LINEAR CORE (19 direct input-$>$output edges, z-scored inputs): ca +1.480, sex:0 -1.072, slope:1 -0.855, chol +0.546, cp:1 -0.493, trestbps +0.480, cp:3 -0.440, exang:1 +0.431, thal:3 -0.410, thal:6 -0.384, cp:4 +0.377, slope:3 -0.366, fbs:0 +0.278, thalach -0.207, thal:4 +0.183, slope:2 +0.153, exang:0 +0.149, age +0.117, thal:7 -0.077. Four inputs are unconnected (sex:1, fbs:1, restecg:1, restecg:2; zero-effect entries of this root). Oldpeak, restecg:0 and cp:2 have no direct edge and act only through the chain. COMPOSITION: serial and additive. The three children form a chain, each read out additively into z: ischaemia\_\allowbreak{}gate (oldpeak, normal resting ECG, reversible defect) feeds bp\_\allowbreak{}offset\_\allowbreak{}gate (the ischaemia score minus a resting-BP offset), which in turn switches off age\_\allowbreak{}gate\_\allowbreak{}without\_\allowbreak{}ischaemia (an extra age term for patients without ischaemia evidence). The chain therefore raises risk with ST depression, more steeply where resting BP is low, and adds a small age term where the chain is quiet. Identity nodes identity\_\allowbreak{}290 and identity\_\allowbreak{}1987 carry the chain links (204 -$>$ 290, 290 -$>$ 1987) so that each gate is its own leaf. EVIDENCE (train, n=242): full model AUC 0.9418, acc 0.872 (compute\_\allowbreak{}performance)\allowbreak{}. Linear core alone AUC 0.934; logit SD 3.14 (core) vs 0.67 (chain). Removing the whole chain (node 204 with downstream) gives 0.934; readout-only ablations: ischaemia 0.940, BP 0.939, age 0.941. Treating all three ReLUs as identities gives 0.936; one least-squares linear fit of the chain on its seven inputs gives 0.937 (R2 0.84); an L2 logistic regression on all 26 inputs gives 0.936. Mean $|$SHAP$|$ (train, all 242): ca 0.144, sex:0 0.110, slope:1 0.100, thal:3 0.048, exang:1 0.047, chol 0.045, cp:4 0.044, cp:3 0.044. Interpretation: close to logistic regression plus one ischaemia correction: ST depression raises risk only above a threshold set by a normal resting ECG and a reversible defect, more so when resting BP is low; age gets a small extra weight when there is no such evidence. Counter-directional signs reported, not explained: restecg:0 positive in the ischaemia gate, trestbps negative in the BP gate, fbs:0 +0.28 directly. Same design as the 10-seed heart-s5 record (generation 37, key 9330: ischaemia gate -$>$ resting-BP gate -$>$ age term), here with three hidden nodes instead of six.

\textbf{Closed form} (children as named functions):
{\small\begingroup\allowdisplaybreaks
\begin{align*}
y_{0} &= \sigma\bigl(z_{0}\bigr) \\
z_{0} &= -0.504 + 0.143 \operatorname{ischaemia\_gate}\bigl(x_{restecg:0}, x_{oldpeak},  \\
&\qquad x_{thal:7\_b}\bigr) + 1.48 x_{ca} \\
&\qquad + 0.143 \operatorname{age\_gate\_without\_ischaemia}\bigl(- 1.59 \operatorname{bp\_offset\_gate}\bigl(1.08 \operatorname{ischaemia\_gate}\bigl(x_{restecg:0}, \\
&\qquad x_{oldpeak}, x_{thal:7\_b}\bigr), x_{cp:2}, x_{cp:3\_b}, x_{trestbps\_b}\bigr), \\
&\qquad x_{age\_b}\bigr) + 0.431 x_{exang:1} + 0.48 x_{trestbps\_a} + 0.546 x_{chol} + 0.278 x_{fbs:0} \\
&\qquad + 0.149 x_{exang:0} + 0.153 x_{slope:2} \\
&\qquad + 0.17 \operatorname{bp\_offset\_gate}\bigl(1.08 \operatorname{ischaemia\_gate}\bigl(x_{restecg:0}, \\
&\qquad x_{oldpeak}, x_{thal:7\_b}\bigr), x_{cp:2}, x_{cp:3\_b}, x_{trestbps\_b}\bigr) \\
&\qquad + 0.183 x_{thal:4} + 0.377 x_{cp:4} + 0.117 x_{age\_a} - 0.207 x_{thalach} - 0.41 x_{thal:3} \\
&\qquad - 0.493 x_{cp:1} - 0.0766 x_{thal:7\_a} - 0.384 x_{thal:6} - 1.07 x_{sex:0} \\
&\qquad - 0.366 x_{slope:3} - 0.855 x_{slope:1} - 0.44 x_{cp:3\_a}
\end{align*}\endgroup}

Expanded over the children:
{\small\begingroup\allowdisplaybreaks
\begin{align*}
y_{0} &= \sigma\bigl(z_{0}\bigr) \\
z_{0} &= -0.504 + 0.143 \operatorname{relu}\bigl(0.247 + 1.15 x_{thal:7\_b} + 2.59 x_{oldpeak} \\
&\qquad + 0.87 x_{restecg:0}\bigr) + 1.48 x_{ca} + 0.143 \operatorname{relu}\bigl(0.203 \\
&\qquad + 0.875 x_{age\_b} - 1.59 \operatorname{relu}\bigl(-0.495 + 0.0164 x_{cp:3\_b} \\
&\qquad + 1.08 \operatorname{relu}\bigl(0.247 + 1.15 x_{thal:7\_b} + 2.59 x_{oldpeak} \\
&\qquad + 0.87 x_{restecg:0}\bigr) + 0.176 x_{cp:2} - 0.815 x_{trestbps\_b}\bigr)\bigr) \\
&\qquad + 0.431 x_{exang:1} + 0.48 x_{trestbps\_a} + 0.546 x_{chol} + 0.278 x_{fbs:0} \\
&\qquad + 0.149 x_{exang:0} + 0.153 x_{slope:2} + 0.17 \operatorname{relu}\bigl(-0.495 \\
&\qquad + 0.0164 x_{cp:3\_b} + 1.08 \operatorname{relu}\bigl(0.247 + 1.15 x_{thal:7\_b} \\
&\qquad + 2.59 x_{oldpeak} + 0.87 x_{restecg:0}\bigr) + 0.176 x_{cp:2} - 0.815 x_{trestbps\_b}\bigr) \\
&\qquad + 0.183 x_{thal:4} + 0.377 x_{cp:4} + 0.117 x_{age\_a} - 0.207 x_{thalach} - 0.41 x_{thal:3} \\
&\qquad - 0.493 x_{cp:1} - 0.0766 x_{thal:7\_a} - 0.384 x_{thal:6} - 1.07 x_{sex:0} \\
&\qquad - 0.366 x_{slope:3} - 0.855 x_{slope:1} - 0.44 x_{cp:3\_a}
\end{align*}\endgroup}

where $y_{0}$ is node \texttt{0}; $z_{i}$, where shown, is the pre-activation of a sigmoid output $y_{i}$; weights and biases are shown to 3 significant figures.
 On the training split (n = 242), the displayed closed form differs from the forward pass by at most 0.00186, and the exact function by at most $2.78\times 10^{-16}$.

\textbf{Evidence} (2 items; full payloads in the JSON record):
\begin{itemize}
\item \textsc{analytical}, \texttt{analytical\_\allowbreak{}formula}; data split: n/a (analytical). Root closed form (get\_\allowbreak{}formula force=true); output is sigmoid.
\item \textsc{empirical}, \texttt{empirical\_\allowbreak{}viz}; data split: train. Train-split performance, ablations, linear replacements and SHAP.
\end{itemize}

\textbf{Authorship:} drafted by AI-assisted draft on 2026-11-10; verified by nobody (unverified).

\subsubsection{\texorpdfstring{A2: ischaemia\_\allowbreak{}gate}{A2: ischaemia gate}}\label{app:heart-s5-c1:ann-ischaemia-gate}
% annotation id ann_32, operation seq 32
\textbf{Kind:} Leaf annotation (depth 1); parent A1 (\cref{app:heart-s5-c1:ann-heart-disease-logit}).

\begin{itemize}
\item \textbf{Entry nodes:} \texttt{-12} (restecg:0), \texttt{-18} (oldpeak), \texttt{-26\_\allowbreak{}b} (thal:7\_\allowbreak{}b).
\item \textbf{Exit nodes:} \texttt{204}.
\item \textbf{Internal nodes:} none.
\item \textbf{Edges (weight):} \texttt{-12}$\to$\texttt{204}~(0.871), \texttt{-18}$\to$\texttt{204}~(2.59), \texttt{-26\_\allowbreak{}b}$\to$\texttt{204}~(1.14).
\item \textbf{Covers} (coverage, own and as part of the hierarchy): 3 candidate nodes.
\end{itemize}

\textbf{Hypothesis.} Ischaemia-evidence gate (raises risk). 204 = ReLU(0.247 + 2.593*oldpeak + 0.871*restecg:0 + 1.145*thal:7). Hypothesis: a one-sided ST-depression score that opens above an oldpeak threshold set by two flags (always open with a normal resting ECG and a reversible thallium defect; above z-oldpeak -0.09 with a normal ECG only, -0.33 with a defect only, +0.58 with neither), and then raises the logit with oldpeak, directly (0.143 per unit) and through the blood-pressure gate (a further 0.183 per unit while that gate is open). EVIDENCE (train, n=242): active for 48\% (100\% of the 44 patients with both flags, 20\% of the 75 with neither); corr with oldpeak 0.88; disease rate 66\% when active vs 24\% when inactive. Mean $|$sensitivity$|$ oldpeak 1.24 $>$ thal:7 0.55 $>$ restecg:0 0.42. This is the main nonlinear effect: removing it with its downstream effects lowers train AUC 0.942 -$>$ 0.934 (to the linear-core level); an affine replacement recovers 0.938. Supported. Oldpeak and restecg:0 have no direct edge; they act only here. The positive restecg:0 weight (normal ECG raises the score) is counter-intuitive and is reported, not explained.

\textbf{Closed form:}
{\small\begingroup\allowdisplaybreaks
\begin{align*}
y_{0} &= \operatorname{relu}\bigl(0.247 + 1.15 x_{thal:7\_b} + 2.59 x_{oldpeak} \\
&\qquad + 0.87 x_{restecg:0}\bigr)
\end{align*}\endgroup}

where $y_{0}$ is node \texttt{204}; $z_{i}$, where shown, is the pre-activation of a sigmoid output $y_{i}$; weights and biases are shown to 3 significant figures.
 On the training split (n = 242), the displayed closed form differs from the forward pass by at most 0.0113, and the exact function by at most $1.78\times 10^{-15}$.

\textbf{Evidence} (2 items; full payloads in the JSON record):
\begin{itemize}
\item \textsc{analytical}, \texttt{analytical\_\allowbreak{}formula}; data split: n/a (analytical); thresholds: thresholds = z-oldpeak at which 204 opens: always open with restecg:0 an.... Closed form from get\_\allowbreak{}formula (force) with readouts and thresholds.
\item \textsc{empirical}, \texttt{empirical\_\allowbreak{}viz} (sensitivity (max\_\allowbreak{}samples 100000)); data split: train. Train-split sensitivity, activity, outcome rates and ablation.
\end{itemize}

\textbf{Authorship:} drafted by AI-assisted draft on 2026-11-10; verified by nobody (unverified).

\subsubsection{\texorpdfstring{A3: bp\_\allowbreak{}offset\_\allowbreak{}gate}{A3: bp offset gate}}\label{app:heart-s5-c1:ann-bp-offset-gate}
% annotation id ann_33, operation seq 33
\textbf{Kind:} Leaf annotation (depth 1); parent A1 (\cref{app:heart-s5-c1:ann-heart-disease-logit}).

\begin{itemize}
\item \textbf{Entry nodes:} \texttt{identity\_\allowbreak{}290}, \texttt{-5} (cp:2), \texttt{-6\_\allowbreak{}b} (cp:3\_\allowbreak{}b), \texttt{-8\_\allowbreak{}b} (trestbps\_\allowbreak{}b).
\item \textbf{Exit nodes:} \texttt{290}.
\item \textbf{Internal nodes:} none.
\item \textbf{Edges (weight):} \texttt{identity\_\allowbreak{}290}$\to$\texttt{290}~(1), \texttt{-5}$\to$\texttt{290}~(0.176), \texttt{-6\_\allowbreak{}b}$\to$\texttt{290}~(0.0164), \texttt{-8\_\allowbreak{}b}$\to$\texttt{290}~(-0.815).
\item \textbf{Covers} (coverage, own and as part of the hierarchy): 4 candidate nodes.
\end{itemize}

\textbf{Hypothesis.} Blood-pressure offset gate on the ischaemia score (raises risk). 290 = ReLU(-0.495 + 1.079*ischaemia\_\allowbreak{}gate - 0.815*trestbps + 0.176*cp:2 + 0.016*cp:3). Hypothesis: passes the ischaemia score on, minus an offset that grows with resting blood pressure; it adds 0.170 per unit to the logit and switches off the age gate. With the ischaemia gate closed it still opens for low resting blood pressure, so within this gate lower resting BP raises risk (0.170*0.815 = 0.139 per SD), partly offsetting trestbps's direct +0.48 edge. cp:2 (atypical angina) nudges it open; cp:3's weight 0.016 is negligible. EVIDENCE (train, n=242): active for 58\% (83\% with the ischaemia gate open, 35\% without); with the ischaemia gate open, active for 100\%, 87\% and 59\% of the low, middle and high resting-BP groups. Disease rate 51\% when active vs 35\% when inactive. Mean $|$sensitivity$|$ ischaemia score 0.58 $>$ trestbps 0.45 $>$ cp:2 0.10 $>$ cp:3 0.01. Readout ablation lowers train AUC 0.942 -$>$ 0.939; an affine replacement fits it well (R2 0.97, AUC 0.941), so its nonlinearity matters little. Supported, small. The negative BP term inside a risk-raising gate is reported, not explained.

\textbf{Closed form:}
{\small\begingroup\allowdisplaybreaks
\begin{align*}
y_{0} &= \operatorname{relu}\bigl(-0.495 + 1.0 x_{identity\_290} + 0.0164 x_{cp:3\_b} \\
&\qquad + 0.176 x_{cp:2} - 0.815 x_{trestbps\_b}\bigr)
\end{align*}\endgroup}

where $y_{0}$ is node \texttt{290}; $z_{i}$, where shown, is the pre-activation of a sigmoid output $y_{i}$; weights and biases are shown to 3 significant figures.
 On the training split (n = 242), the displayed closed form differs from the forward pass by at most $8.86\times 10^{-4}$, and the exact function by at most $8.88\times 10^{-16}$.

\textbf{Evidence} (2 items; full payloads in the JSON record):
\begin{itemize}
\item \textsc{analytical}, \texttt{analytical\_\allowbreak{}formula}; data split: n/a (analytical). Closed form from get\_\allowbreak{}formula (force) with readouts and thresholds.
\item \textsc{empirical}, \texttt{empirical\_\allowbreak{}viz} (sensitivity (max\_\allowbreak{}samples 100000)); data split: train. Train-split sensitivity, activity, outcome rates and ablation.
\end{itemize}

\textbf{Authorship:} drafted by AI-assisted draft on 2026-11-10; verified by nobody (unverified).

\subsubsection{\texorpdfstring{A4: age\_\allowbreak{}gate\_\allowbreak{}without\_\allowbreak{}ischaemia}{A4: age gate without ischaemia}}\label{app:heart-s5-c1:ann-age-gate-without-ischaemia}
% annotation id ann_34, operation seq 34
\textbf{Kind:} Leaf annotation (depth 1); parent A1 (\cref{app:heart-s5-c1:ann-heart-disease-logit}).

\begin{itemize}
\item \textbf{Entry nodes:} \texttt{identity\_\allowbreak{}1987}, \texttt{-1\_\allowbreak{}b} (age\_\allowbreak{}b).
\item \textbf{Exit nodes:} \texttt{1987}.
\item \textbf{Internal nodes:} none.
\item \textbf{Edges (weight):} \texttt{identity\_\allowbreak{}1987}$\to$\texttt{1987}~(1), \texttt{-1\_\allowbreak{}b}$\to$\texttt{1987}~(0.875).
\item \textbf{Covers} (coverage, own and as part of the hierarchy): 2 candidate nodes.
\end{itemize}

\textbf{Hypothesis.} Age gate without ischaemia evidence (raises risk with age, small). 1987 = ReLU(0.203 + 0.875*age - 1.589*\allowbreak{}bp\_\allowbreak{}offset\_\allowbreak{}gate)\allowbreak{},\allowbreak{} read out +0.143. Hypothesis: an extra age term (0.125 per SD of age, on top of age's direct +0.117) that acts only while the BP gate is closed or nearly closed, i.e. for patients without much ischaemia evidence; any opening of the BP gate cancels it (each unit of the BP gate removes 0.227 from the logit through this node until it closes). EVIDENCE (train, n=242): active for 35\% (65\% when the BP gate is closed, rising with age: 0\%, 80\%, 100\% of the youngest, middle and oldest groups; 14\% when the BP gate is open). Disease rate barely differs (43\% vs 45\%). Mean $|$sensitivity$|$ BP gate 0.35, age 0.31. Removing it lowers train AUC only 0.942 -$>$ 0.941. Supported, near-negligible.

\textbf{Closed form:}
{\small\begingroup\allowdisplaybreaks
\begin{align*}
y_{0} &= \operatorname{relu}\bigl(0.203 + 1.0 x_{identity\_1987} + 0.875 x_{age\_b}\bigr)
\end{align*}\endgroup}

where $y_{0}$ is node \texttt{1987}; $z_{i}$, where shown, is the pre-activation of a sigmoid output $y_{i}$; weights and biases are shown to 3 significant figures.
 On the training split (n = 242), the displayed closed form differs from the forward pass by at most $4.99\times 10^{-4}$, and the exact function by at most $2.22\times 10^{-16}$.

\textbf{Evidence} (2 items; full payloads in the JSON record):
\begin{itemize}
\item \textsc{analytical}, \texttt{analytical\_\allowbreak{}formula}; data split: n/a (analytical). Closed form from get\_\allowbreak{}formula (force) with readouts and thresholds.
\item \textsc{empirical}, \texttt{empirical\_\allowbreak{}viz} (sensitivity (max\_\allowbreak{}samples 100000)); data split: train. Train-split sensitivity, activity, outcome rates and ablation.
\end{itemize}

\textbf{Authorship:} drafted by AI-assisted draft on 2026-11-10; verified by nobody (unverified).

\subsection{Audit against the definitions}\label{app:heart-s5-c1:audit}
Automatic findings on the stored explanation (codes as in the JSON record's \texttt{issues}); they flag what an examiner should check, not verdicts.
\begin{itemize}
\item \textsc{warning} \texttt{unverified}: No human verification is recorded for this explanation.
\end{itemize}
\end{drafted}
